% Generated draft: requires mathematical and source review.
\documentclass[11pt]{article}
\usepackage[margin=25mm]{geometry}
\usepackage{fontspec}
\setmainfont{Latin Modern Roman}
\usepackage{amsmath,amssymb,mathtools,unicode-math}
\setmathfont{Latin Modern Math}
\usepackage{enumitem,needspace,xurl,hyperref}
\setlength{\parindent}{0pt}
\setlength{\parskip}{5pt}
\setlength{\emergencystretch}{3em}
\newcommand{\N}{\mathbb{N}}
\newcommand{\Z}{\mathbb{Z}}
\newcommand{\Q}{\mathbb{Q}}
\newcommand{\R}{\mathbb{R}}
\newcommand{\C}{\mathbb{C}}
\newcommand{\F}{\mathbb{F}}
\newcommand{\A}{\mathbb{A}}
\newcommand{\PP}{\mathbb{P}}
\newcommand{\E}{\mathbb{E}}
\newcommand{\eps}{\varepsilon}
\newcommand{\dd}{\,\mathrm{d}}
\begin{document}
\section*{UnsolvedMath ID: 701075}
\textbf{Source problem code:} catalogue entry
\subsection{Definitions and mathematical statement}
Let \(p\) be a fixed prime. A finite \(p\)-group is a finite group \(G\) with \(|G|=p^n\) for some integer \(n\geq 0\).

The supplied question asks whether, for every \(a\geq 0\), there exists a finite \(p\)-group of abundance \(a\). Schematically, this is
\[
\forall a\geq 0\quad \exists\,G\text{ a finite }p\text{-group}\quad \operatorname{abundance}(G)=a.
\]
Here \(\operatorname{abundance}(G)\) is only a placeholder for the source's undefined term ``abundance'', not a defined invariant. The supplied record neither defines abundance nor specifies the domain of \(a\). Consequently, this display records the requested quantifier structure but is not a fully specified mathematical statement. Both the definition of abundance and the intended range of \(a\) must be supplied before the question can be made precise.
\subsection{Short English statement}
For each fixed prime \(p\), the question asks whether every nonnegative admissible value \(a\) occurs as the abundance of a finite \(p\)-group. The supplied record leaves both abundance and the domain of \(a\) undefined.
\subsection{Sources}
\begin{itemize}
\item Supplied local UnsolvedMath registry record: permanent ID \(701075\), exact source problem code \texttt{catalogue entry}, titled ``Kourovka Notebook Problem 14.91''.
\item The record attributes the problem to G. Fern\'andez--Alcober, proposed in 1999, and cites \emph{The Kourovka Notebook (Unsolved Problems in Group Theory)}, 14th issue (1999), Problem 14.91; it also identifies the 21st issue edition (2026), \texttt{arXiv:1401.0300v48}.
\item The supplied collection description identifies E. I. Khukhro and V. D. Mazurov as editors. The registry labels the problem ``open''; this is a supplied classification, not an independently verified mathematical or literature assessment. The cited source has not been independently consulted for this draft.
\end{itemize}
\end{document}
